\documentclass{article}
\usepackage{amsmath}
\usepackage{caption}
\usepackage[UTF8]{ctex}

\begin{document}
\title{Tool Report}
\author{曾奕豪 - 25100021005}
\maketitle

\section{Result}

% 正文引用公式：\eqref{标签名} 会自动生成带括号的公式编号
如式 \eqref{eq:sum} 所示，这是一个标准的自然数求和公式。
% equation 环境：自动生成公式编号
\begin{equation}
    \sum_{i=1}^{n} i = \frac{n(n+1)}{2}
    \label{eq:sum}  % \label 给公式打标签，用于交叉引用
\end{equation}

% 正文引用表格：\ref{标签名} 会自动生成表格编号
实验的核心参数统计如表 \ref{tab:params} 所示。
% table 浮动体环境：存放表格、标题和标签
\begin{table}[h]
    \centering % 表格居中
    \caption{实验核心参数表} % 表格标题
    \label{tab:params} % 表格标签，用于引用
    % 2行2列表格：{|c|c|} 表示两列居中、带竖线
    \begin{tabular}{|c|c|}
        \hline
        参数名称 & 参数数值 \\
        \hline
        平均延迟 & 193.75 ms \\
        \hline
    \end{tabular}
\end{table}
\end{document}
